% Budget: 10 pp.  Function: what we found.  Every number carries \prov{}{}{}.
% ORDER MATTERS: the positive result comes BEFORE the negative one.  A
% positive cell + a matched architectural negative + the causal variable
% that discriminates them is a FINDING; the reverse order reads as a
% post-mortem with a footnote.  See rubric/01-originality.md.
\chapter{Results}
\label{ch:results}

\section{The framework across backbones (D1)}
\label{sec:res-d1}
\todo{Three backbone families behind one interface + the port of AVID's own
recipe into AVID's own repository on a new base family. Plus the
LoRA-vs-output comparison once it lands.}

\section{A working action-conditioned cell (D2)}
\label{sec:res-working-cell}
\todo{THE POSITIVE RESULT. DC x ACWM Robot Arm, three independent readouts:
effect_rel at 3.9x the AVID reference; all three structure probes above
chance on held-out ind_test; and the action-following solution has a LOWER
denoising loss than the untreated control -- it is not a trade against
prediction quality. The untreated control clears the reference unaided, so
the cell works natively and condition_center is a ~6x accelerator, not an
enabler. This is also the ablation's positive control.

LIMITS THAT TRAVEL WITH THE CLAIM: runs cancelled pre-convergence, so quote
the acceleration and never a level; rollout-level control is NOT measured on
this cell (the swap probe exists only on the Wan path).}

\subsection{Output quality on the working cell (qualitative)}
\label{sec:res-dc-quality}
\fig{DC rollouts, adapted vs frozen base -- arm E \code{6oyu1inq}, arm F
\code{86kb01su}. Fixed seed, identical initial frame, held-out clips.
Provenance \nv{artefact path + step}.}
\todo{This is the ONLY evidence about the DC cell's output quality: no DC
run logs quality metrics (all 18 checked in dc-acwm-robotarm-avid-parity),
so the gap is otherwise open on the cell that carries the positive result.

STRICT SCOPE -- label qualitative, claim nothing about actions. An
adapted-vs-base comparison cannot separate action-following from a better
temporal prior, and our own Wan cell is the counterexample: it beats its
frozen base on 6/6 quality metrics while every structure probe sits at
chance. The action claim in this section rests on the probes above, NOT on
these videos.}

\section{The injection pathway decides (D2)}
\label{sec:res-pathway}

On the base chosen in \cref{sec:backbones} for its capability, horizon and
sampling cost, the adapter did not use the actions. This section establishes
what caused that, by eliminating the alternatives in turn and then varying
the surviving one directly.

\subsection{The failure}

The cell improves the video substantially. Measured against its own frozen
base on held-out clips it is better on every quality metric we
record.\prov{vy9tcuco}{outputs/acwm-robotarm-wan-tokennorm-nobase/checkpoints/}{\nv{commit}}
On the same checkpoint, all three structure probes of \cref{sec:probes}
return chance: the steering cosine is $0.00$ against a chance level of
$0.000$, and temporal alignment and spatial concentration are
indistinguishable from theirs. The frozen-base null is exactly zero, so the
probes are measuring the adapter.

An adapter that improves generation while carrying no recoverable action
information is performing domain adaptation, not conditioning. The remainder
of this section asks why.

\subsection{What it is not}

Four explanations are available before reaching the architecture, and each
is excluded by a measurement rather than by argument.

\paragraph{Not the data.} The unmodified published recipe of
\citet{rigter2024avid} follows actions on the same environment, with
the same frozen weights and the same probe, where three of our adapters did
not.\prov{rqp4s3gp}{\nv{ckpt}}{\nv{commit}} Whatever prevents the action
being used is on our side of the experiment.

\paragraph{Not the adapter's capacity.} A structurally clean adapter of
\SI{7.5}{\mega\nothing} parameters settles \emph{below} the larger
DiT-derived arms rather than above them, so the constraint is not that the
module is too small to represent the mapping.

\paragraph{Not the family.} The same adapter family, on a different frozen
base, reduces the denoising loss by \SI{74.9}{\percent} and contributes
\num{0.52} of the prediction, against \SI{3.3}{\percent} and \num{0.047} on
the comparable arm of the base under
study.\prov{25240257}{\nv{ckpt}}{\nv{commit}} On that base the adapted
prediction is numerically almost the frozen one, its cosine to the base
being \num{0.9989}. Output adapters can reshape a frozen video model
substantially; this one did not.

\paragraph{Not a merely mis-scaled or mis-aligned signal.} Rescaling the
action pathway and aligning the action tokens to the latent grid both
change the magnitude of the response without changing its structure, and
substituting channel concatenation for cross-attention performs no better at
any matched step.\prov{6ruz55f6}{outputs/acwm-robotarm-CONCAT-nobase-run/checkpoints/}{\nv{commit}}

\subsection{Where the signal is lost}

Tracing the action through the network locates the failure precisely. The
action tokens are informative on entry and remain so throughout the
attention stack: the cross-attention output is between \SI{44}{\percent} and
\SI{56}{\percent} action-driven in every one of the ten conditioned blocks.
The relative sensitivity of the stream to the action then falls from
\num{0.44} to \num{0.0085} across a \emph{single} operation, the residual
addition.\prov{25085110}{ncztxyyo}{\nv{commit}}

The magnitudes involved make the arithmetic unsurprising once measured. The
cross-attention output has a root-mean-square of approximately \num{0.01}
against a residual stream of \numrange{1.8}{3.0}. The signal survives the
attention and is then added to something two orders of magnitude larger.

\subsection{Varying the pathway directly}

The preceding is a diagnosis, and a diagnosis invites the objection that
something else differed. We therefore compare two injection routes with
everything else held fixed: the same frozen base, the same data, and, at the
point of comparison, matched adapter contribution (\num{0.114} against
\num{0.111}) and matched composition mask (\num{0.953}). The only difference
is whether the per-frame action is supplied as cross-attention tokens or as
a per-frame modulation of normalised activations.

The modulated arm carries more action information at every depth measured,
and the gap widens with training rather than closing: a factor of
\num{2.49} at step \num{12000}, Welch $t = 10.5$, with the two arms'
contributions identical to within \SI{0.2}{\percent}.\prov{avid\_wan\_rt1\_47M\_faithful}{\nv{ckpt}}{\nv{commit}}
The action-driven share of the modulated arm rises from \SI{15.3}{\percent}
to \SI{25.3}{\percent} over that interval while the token arm moves from
\SI{9.1}{\percent} to \SI{10.2}{\percent}, which is the signature of an
information ceiling rather than of under-training.

\paragraph{The measurement is not a gain artefact.} Because
\cref{eq:effect-rel} is monotone in the gain of the action pathway
(\cref{sec:integrity}), a larger response is by itself compatible with a
louder signal rather than a better-used one. The temporal control separates
them: the modulated arm concentrates its per-frame response on the diagonal
at \num{0.409} against a chance level of \num{0.200}, while the token arm
sits at \num{0.198}, and its per-frame response rows are \emph{bit-identical}
to one another. Gain scales a response; it cannot manufacture per-frame
differentiation where none exists.

\subsection{Scope of the claim}

Two boundaries matter, because the unscoped version is refutable.

The conditioning signal here is \emph{low-dimensional and spatially global}:
a seven-degree-of-freedom vector with no spatial extent. For signals that
vary across the frame, the ordering reverses, and cross-attention's
finer-grained spatial treatment is what recommends
it~\citep{chen2024gentron}; the same has been reported for
high-degree-of-freedom dexterous conditioning. Our result is a statement
about the regime in which a global modulation is expressive enough to carry
the signal.

The base is also \emph{frozen} and the adapter is trained from scratch
beside it, which is what makes the magnitude comparison severe. In models
trained end-to-end the cross-attention pathway operates at a substantially
higher fraction of the residual stream and is demonstrably not
inert~\citep{motif2026video}. The claim is therefore conditional: in a
frozen backbone with a from-scratch adapter carrying a low-dimensional
signal, the cross-attention write lands roughly an order of magnitude below
the level at which jointly trained models operate that pathway.

\paragraph{What is ours and what is not.} That modulation of normalised
activations outperforms token-based conditioning for low-dimensional signals
is established, in class-conditional
diffusion~\citep{peebles2023dit} and, more recently, for action
conditioning specifically~\citep{bai2026ioi, huang2026nano}; that the
gap widens under a frozen backbone has also been
reported~\citep{zhou2026decoupled}, attributed there to conditioning being routed
through frozen projections. What this section adds is the localisation. The
signal is not lost in the attention computation, where it remains between
\SI{44}{\percent} and \SI{56}{\percent} action-driven; it is lost at the
residual addition. That distinguishes a magnitude account from a
representational one, and it is testable: an account based on stale
projections does not predict that additive injection, which bypasses those
projections entirely, degrades under freezing as well.

\section{Two scale failures at the same interface (D2)}
\label{sec:res-scale}
\todo{DC's learned pedestal; Wan's drowning at the residual add. Both are
scale calibration, in opposite directions, at the same interface.}

\section{What the adapter extracts from the action}
\label{sec:res-decomposition}

The preceding sections establish that an action signal reaches the adapter
and changes its output. They do not establish \emph{what} the adapter takes
from it. This section answers that, and the answer explains two
observations that would otherwise have to be reported as unexplained
negatives.

\subsection{Three measurements}

\paragraph{The response has no direction.} On the cross-attention cell the
structure probes of \cref{sec:probes} return chance on all three axes: the
steering cosine is $0.00$ against a chance level of $0.000$, so the
direction in which an action moves the prediction is uncorrelated with the
direction the data implies; temporal alignment and spatial concentration are
likewise indistinguishable from chance, so the response is neither placed in
the frames the action governs nor localised on the parts of the scene that
move.\prov{probe job 25107536}{acwm-robotarm-wan-tokennorm-nobase}{\nv{commit}}
The signal is present and measurable, and it is structureless.

\paragraph{The response does carry magnitude.} On the distilled base a
paired control isolates the action directly: identical weights, identical
conditioning frame, identical seed, with only another clip's actions
substituted. Per-clip motion in the adapted rollout tracks ground-truth
motion better than in the shuffled-action rollout in \textbf{four of four
draws}: gains of $+0.069$, $+0.168$, $+0.122$ and $+0.213$, mean
$+0.143$, sign test $p \approx 0.06$ one-tailed.\prov{mo3k2639}{outputs/wan22turbo-action-robotarm/checkpoints/}{\nv{commit}}
Because the two rollouts differ in exactly one respect, this is a causal
quantity rather than a correlational one.

\paragraph{The response does not carry correctness.} Across every cell we
measure, supplying the \emph{correct} action does not reduce prediction
error relative to supplying a wrong one. The action loss gap finishes at
$0.00005$ on the distilled arm and remains within $|{\cdot}| \leq 0.00055$
across all ten evaluations of its predecessor, with the action cosine never
leaving $0.9998$.

\subsection{The reading}

Taken together these are not three separate findings:

\begin{quote}
\emph{The adapter reads the action as a scalar (roughly, ``how much
movement'') rather than as a directional command.}
\end{quote}

A scalar response accounts for all three measurements at once, and it
accounts for them in the right direction rather than merely being consistent
with them.

It explains the structure probes. A quantity with no direction has no
steering direction to detect, no frame to be localised to and no spatial
locus to concentrate on; the triad returning chance is then a correct
measurement of a scalar, not a failure of the instrument or a null result
about the adapter.

It explains the loss gap. Squared error penalises motion of the right
magnitude in a slightly wrong place as heavily as no motion at all, so a
model that has learned how much the arm should move, but not where,
gains nothing on the objective. The magnitude information is real, and the
metric cannot see it.

And it sharpens the economics of \cref{sec:res-scale}. Under a
teacher-forced denoising objective the actions account for a very small
fraction of the loss; what the gradient does buy with that fraction is the
cheapest useful summary of the action available, which is its magnitude.
A directional response would cost more and pay no better.

\subsection{Scope}

The magnitude effect is modest and its interval is not yet established: the
confidence intervals available are on the correlation with ground truth
rather than on the gain, and that correlation is itself indistinguishable
from zero in three of the four draws. The claim rests on the sign of the
gain, which is consistent across all four, from a single run without seed
replication, on one base. We therefore state it as \emph{a modest,
consistently-signed effect on motion magnitude} and not as trajectory
control.

That distinction is the boundary of what this thesis demonstrates. An
adapter on a frozen video prior can be made to register the action and to
modulate the amount of motion accordingly; making it produce the motion the
action specifies is a different problem, and \cref{sec:boundary} argues that
it is the objective rather than the architecture that stands in the way.

\section{Standard metrics are blind to whether the conditioning is used}
\label{sec:res-blindness}

The preceding sections separated two behaviours that a conventional
evaluation reports identically. We now state that observation directly,
because it is the finding with the widest reach beyond this thesis: \emph{no
standard readout used to monitor a conditional generative model can
distinguish a model that uses its conditioning from one that does not.}

\subsection{The demonstration}

Consider the Wan~$\times$~ACWM cell without base-prediction injection
\provshort{vy9tcuco}. Measured against its own frozen base on held-out
clips, the adapted model is better on \emph{every} quality metric we
record: FID improves by \SI{27.5}{\percent}, FVD-I3D by
\SI{45.1}{\percent}, LPIPS by \SI{16.0}{\percent}, and PSNR, MSE and SSIM by
\num{9.7}, \num{31.0} and \SI{1.7}{\percent} respectively.\prov{vy9tcuco}{outputs/acwm-robotarm-wan-tokennorm-nobase/checkpoints/}{\nv{commit}}
By any measure a practitioner would normally consult, the adapter works.
With base-prediction injection enabled the same picture is stronger still,
FVD improving by \SI{63.7}{\percent}, from \num{1118.4} to \num{406.3}
\provshort{52o3uxz8}.

On that same checkpoint, all three structure probes of
\cref{sec:probes} return chance. The steering cosine is \num{0.00} against a
chance level of \num{0.000}: the direction in which an action moves the
prediction is uncorrelated with the direction the data implies. Temporal
alignment and spatial concentration are likewise indistinguishable from
their chance levels, so the response is neither placed in the frames the
action governs nor localised on the parts of the scene that move.\prov{probe job 25107536}{acwm-robotarm-wan-tokennorm-nobase}{\nv{commit}}
The frozen-base null is exactly zero throughout, so the probes are measuring
the adapter and not a leak.

A \SI{45}{\percent} improvement in FVD that carries no recoverable action
information is a \emph{domain} correction. The adapter has learned the
distribution of this environment's videos, which is a genuine and useful
thing to learn, and it is not the thing we asked it to learn.

The complementary case makes the same point from the other side. The
DynamiCrafter cell of \cref{sec:res-working-cell} does follow actions (all
three structure probes clear chance on held-out data), and \emph{no run in
that campaign logs a single quality metric}. Across eighteen runs there is
no FID, no FVD, and no LPIPS. The cell that demonstrably uses its actions is
the one for which we cannot report the numbers a reader would expect, and
the cell with the strongest quality numbers is the one that ignores them.
Had we monitored either cell in the conventional way, we would have drawn
the opposite conclusion about both.

\subsection{Why each readout fails}

The blindness is not an accident of these particular runs; each readout is
insensitive by construction.

\emph{Training loss} is dominated by appearance. On this data actions
explain approximately \SI{0.45}{\percent} of the teacher-forced denoising
objective \provshort{25097452}, so a model that ignored actions entirely
would forfeit less than one percent of the quantity being minimised,
comfortably inside the noise of a training curve.

\emph{Gate and contribution statistics} measure how far the adapter moves
the output from the frozen base, not what it moves it toward. A saturated
gate and a healthy one are both compatible with an action-blind residual.

\emph{Perceptual metrics} compare distributions of generated and real video.
An adapter that sharpens the environment's visual statistics improves them;
so does one that also respects the actions. FID and FVD have no term that
depends on the conditioning at all.

\emph{Pixel metrics} are worse than uninformative: they reward the
mean-regression behaviour that action-blindness produces. The RT-1 cells
improve on PSNR and MSE while degrading on FID and LPIPS
\provshort{0fqjrqjl}: the signature of a model hedging toward the
conditional mean.

\emph{Visual inspection} fails for the reason the others do. A reader
comparing adapted samples against the frozen base sees a real improvement
and has no access to the counterfactual, namely what this model would have
produced under a different action, which is the only comparison that
carries the answer.

\subsection{What this cost, and what it implies}

The diagnosis required instruments built for the purpose
(\cref{sec:probes}), and building them was not sufficient on its own. Our
first such instrument, the relative action effect, is monotone in the gain
of the action pathway: any intervention that amplifies that pathway raises
the metric whether or not it improves the information the adapter extracts.
Two interventions we had characterised as mechanism repairs are themselves
gain increases, and for a period the distinction was unavailable to us.
Separating them required a probe whose null is a \emph{shape} rather than a
magnitude, and the resolution came from a temporal control in which the
pooled arm's per-frame response rows proved bit-identical, a degeneracy
that no amount of gain can produce (\cref{sec:integrity}).

The generalisable statement is therefore narrower and more useful than
``standard metrics are insufficient''. A scalar summary of a conditional
model's sensitivity cannot separate \emph{more signal} from \emph{louder
signal}; only a measurement with internal structure can. Practitioners
adapting a frozen generative model to a new conditioning variable, and
monitoring the adaptation through loss, gate statistics and sample quality,
have no instrument in that set capable of reporting whether the conditioning
is used; and, on the evidence here, those readouts will look their best
in precisely the case where it is not.

\section{Shortcut on flow versus diffusion (D3)}
\label{sec:res-shortcut}
\todo{The first clean D3 positive. State the cross-base confound IN THE SAME
PARAGRAPH as the 9x, or the number takes the claim with it.}
